% Book 4, Year 2, chapter 6 -- Rust: ownership and borrowing.
% Provenance (never printed): sources/book4/BRIEFS.md, chapter 6.
\chapter{Rust: Ownership and Borrowing}\label{ch:b2:rust-ownership}

Last year's most frequent error, in five lines of C: a block is freed through one pointer and read through another.
The compiler accepts it; only the sanitizer, at run time, on the input that reaches the line, notices.

\omcode{code/bachelor-2/06-rust-ownership/listings/use_after_free.c}{9}{13}{Freed through one pointer, read through
another}\label{lst:b2:rust-ownership:uaf-c}

\omoutput[1-5,7-7]{code/bachelor-2/06-rust-ownership/outputs/run_use_after_free.txt}{A use after free, found by the
sanitizer while the program runs.}\label{out:b2:rust-ownership:uaf-c}

The same program in Rust, the language of this part of the book:

\omcode{code/bachelor-2/06-rust-ownership/listings/use_after_free.rs}{3}{8}{The same program in
Rust}\label{lst:b2:rust-ownership:uaf-rs}

\omoutput[1-12,22-22]{code/bachelor-2/06-rust-ownership/outputs/use_after_free.txt}{Refused by the compiler: no
program is produced.}\label{out:b2:rust-ownership:uaf-rs}

The compiler has found, without running anything, that the string is given back while a reference to it is still
to be used, and it names the three lines involved. It did so by applying a few rules about who owns each value and
who may use it when. This chapter states them; they forbid the use after free, the double free and, as
\cref{ch:b2:threads-and-synchronisation} will show, the data race, in every program that compiles.

\begin{recall}
Book 3: a block of the heap is obtained with \code{malloc} and given back with \code{free}; using it afterwards
(use after free) or freeing it twice (double free) is undefined behaviour, which the address sanitizer detects at
run time. \cref{ch:b2:cpp-classes-raii,ch:b2:cpp-copy-move}: in C++, an object that owns a resource frees it in its
\omterm{def:b2:cpp-classes-raii:constructor}{destructor} (RAII), and a move transfers the resource, leaving the source empty but still usable.
\end{recall}

\section{Rust in a few pages}\label{sec:b2:rust-ownership:basics}

Rust is a compiled language, like C and C++: \code{rustc file.rs} produces an executable, with no garbage collector
and no run-time system beyond C's. Its syntax will look familiar.

\omcode{code/bachelor-2/06-rust-ownership/listings/hello.rs}{2}{19}{Bindings, types, functions, a loop and a
vector}\label{lst:b2:rust-ownership:hello}

\omoutput{code/bachelor-2/06-rust-ownership/outputs/hello.txt}{The fines of four loans, and integer overflow made
visible.}\label{out:b2:rust-ownership:hello}

A few points differ from C. \code{let} introduces a \emph{binding}, immutable unless declared \code{mut}. Types are
written after the name (\code{total: u32}) and are inferred where possible: \code{u8} to \code{u64} and \code{i8} to
\code{i64} are the integers with their sizes in bits, \code{usize} the size of a pointer, \code{f64} the
\code{double}. There is no implicit conversion between numeric types: \code{as} converts. A block is an expression
whose value is its last expression, written without a semicolon, which is how \code{fine_cents} returns. An
arithmetic overflow stops the program in a debug build; \code{checked_add} returns \code{None} instead of
overflowing and \code{wrapping_add} wraps modulo $2^8$ on purpose. \code{println!} is a \emph{macro}, hence the
\code{!}, whose \code{{}} are replaced by the arguments or by the variables named inside them. Tests live in a
\emph{crate}, Rust's unit of compilation, managed by the tool \code{cargo}; the tests of this chapter form the crate
\code{tests/ownership}. Rust 2021 is the edition of the language used throughout, the one that came with Rust
1.56 in October 2021.
% ledger: b2:rust-2021

\section{Ownership}\label{sec:b2:rust-ownership:ownership}

\begin{definition}[Owner, ownership]\label{def:b2:rust-ownership:ownership}
In Rust, every value has exactly one \emph{owner}\index{owner (Rust)}, the variable (or the field, or the element of
a collection) that holds it. When the owner goes out of scope, the value is dropped. \emph{Ownership}\index{ownership}
is this relation, and the rule that it is unique at every moment.
\end{definition}

\begin{definition}[Drop]\label{def:b2:rust-ownership:drop}
To \emph{drop}\index{drop (Rust)} a value is to run its clean-up code, its \code{Drop} implementation if it has one,
then to drop its fields, and finally to free its memory. The compiler inserts the drops, at the end of the \omterm{def:b2:rust-ownership:ownership}{owner}'s
scope, in the reverse order of declaration; \code{drop(x)} drops \code{x} earlier.
\end{definition}

A \code{String} is three words on the stack, a pointer, a length and a capacity, and a buffer on the heap, which it
owns. Assigning it to another variable does not copy the buffer.

\begin{definition}[Move]\label{def:b2:rust-ownership:move}
A \emph{move}\index{move (Rust)} transfers \omterm{def:b2:rust-ownership:ownership}{ownership} of a value: an assignment \code{let b = a;}, the passing of
\code{a} as an argument by value, or its return from a function copies the bytes of \code{a}'s stack part to the
destination, which becomes the \omterm{def:b2:rust-ownership:ownership}{owner}. The source is then \emph{moved from}: the compiler refuses any later use of
it, and does not \omterm{def:b2:rust-ownership:drop}{drop} it.
\end{definition}

\omcode{code/bachelor-2/06-rust-ownership/listings/move_string.rs}{2}{19}{A \omterm{def:b2:rust-ownership:move}{move}, a copy and a
clone}\label{lst:b2:rust-ownership:move}

\omoutput{code/bachelor-2/06-rust-ownership/outputs/move_string.txt}{After \code{clone}, two
buffers.}\label{out:b2:rust-ownership:move}

\begin{omfigure}
\begin{tikzpicture}[font=\scriptsize]
\begin{scope}
\node[font=\footnotesize\bfseries, anchor=west] at (-0.2,1.5) {before \texttt{let b = a;}};
\node[cell, minimum width=12mm] (ap) at (0.8,0.6) {ptr};
\node[cell, minimum width=12mm, below=0mm of ap] (al) {len 11};
\node[cell, minimum width=12mm, below=0mm of al] (ac) {cap 11};
\node[left=1mm of al, font=\scriptsize\ttfamily] {a};
\node[cell, minimum width=30mm, fill=omDef!10, font=\scriptsize\ttfamily] (buf) at (3.6,-2.0) {Quay Street};
\draw[ptr] (ap.east) to[out=0, in=90] (buf.north);
\node[font=\scriptsize, text=black!60] at (0.8,-2.0) {stack};
\node[font=\scriptsize, text=black!60] at (3.6,-2.65) {heap};
\end{scope}
\begin{scope}[xshift=7.0cm]
\node[font=\footnotesize\bfseries, anchor=west] at (-0.2,1.5) {after};
\node[cell, minimum width=12mm, fill=black!8, text=black!45] (ap) at (0.8,0.6) {ptr};
\node[cell, minimum width=12mm, fill=black!8, text=black!45, below=0mm of ap] (al) {len 11};
\node[cell, minimum width=12mm, fill=black!8, text=black!45, below=0mm of al] (ac) {cap 11};
\node[left=1mm of al, font=\scriptsize\ttfamily, text=black!45] {a};
\node[cell, minimum width=12mm] (bp) at (3.0,0.6) {ptr};
\node[cell, minimum width=12mm, below=0mm of bp] (bl) {len 11};
\node[cell, minimum width=12mm, below=0mm of bl] (bc) {cap 11};
\node[left=1mm of bl, font=\scriptsize\ttfamily] {b};
\node[cell, minimum width=30mm, fill=omDef!10, font=\scriptsize\ttfamily] (buf) at (5.0,-2.0) {Quay Street};
\draw[ptr] (bp.east) to[out=0, in=90] (buf.north);
\node[align=center, font=\scriptsize, text=black!60] at (0.8,-2.1) {moved from: unusable,\\ not dropped};
\end{scope}
\end{tikzpicture}
\par\smallskip{\small A \omterm{def:b2:rust-ownership:move}{move} of a \code{String}. The three words of \code{a} are copied into \code{b}; the buffer
stays where it is, with one \omterm{def:b2:rust-ownership:ownership}{owner}, \code{b}. The compiler forbids any use of \code{a} afterwards and \omterm{def:b2:rust-ownership:drop}{drops} only
\code{b}, so the buffer is freed exactly once. In C++, a move would also have emptied \code{a}, and run its
\omterm{def:b2:cpp-classes-raii:constructor}{destructor} later; in Rust nothing runs on \code{a} at all.}
\end{omfigure}

\omoutput[2-10,20-20]{code/bachelor-2/06-rust-ownership/outputs/use_after_move.txt}{Using \code{a} after
\code{let b = a;} is refused.}\label{out:b2:rust-ownership:use-after-move}

\begin{definition}[Copy type]\label{def:b2:rust-ownership:copy-type}
A \emph{copy type}\index{copy type (Rust)} is a type whose values are duplicated by copying their bytes, with no
resource to free, and that implements the \code{Copy} trait: the integers, \code{f64}, \code{bool}, \code{char},
shared references, and tuples and arrays of copy types. For a copy type, assignment and passing by value copy, and
the source stays usable. A type that owns a resource, such as \code{String} or \code{Vec}, cannot be \code{Copy}.
\end{definition}

A deliberate copy of a \code{String} is a \code{clone}: a new buffer with the same bytes. Clones are visible in the
code, which is the point: a C++ copy runs silently at an initialisation or an assignment (\omterm{def:b2:cpp-copy-move:copy}{copy constructor}, copy assignment), a Rust clone
is written.

\section{Borrowing}\label{sec:b2:rust-ownership:borrowing}

Moving every value into every function that reads it would be impractical. A function can instead
\emph{\omterm{def:b2:rust-ownership:borrow}{borrow}} it.

\begin{definition}[Borrow]\label{def:b2:rust-ownership:borrow}
A \emph{borrow}\index{borrow (Rust)} is a reference to a value that does not take \omterm{def:b2:rust-ownership:ownership}{ownership} of it. A \emph{shared
borrow} \code{&x}, of type \code{&T}, allows reading the value; a \emph{mutable}, or \emph{exclusive}, borrow
\code{&mut x}, of type \code{&mut T}, allows changing it. The \emph{borrowing rule}: at any point of the program, a
value has either any number of shared borrows in use, or exactly one exclusive borrow in use, but not both; and no
borrow is used after the value is moved or dropped.
\end{definition}

A \omterm{def:b2:rust-ownership:borrow}{borrow} is \emph{in use} from its creation to its last use, not to the end of its block: the compiler follows the
program's paths, line by line.

\omcode{code/bachelor-2/06-rust-ownership/listings/borrow_rules.rs}{2}{23}{Shared \omterm{def:b2:rust-ownership:borrow}{borrows} to read, an exclusive one to
change}\label{lst:b2:rust-ownership:borrow-rules}

\omoutput{code/bachelor-2/06-rust-ownership/outputs/borrow_rules.txt}{Every \omterm{def:b2:rust-ownership:borrow}{borrow} ended before the next exclusive
one began.}\label{out:b2:rust-ownership:borrow-rules}

\omoutput[2-10]{code/bachelor-2/06-rust-ownership/outputs/two_mut.txt}{Two exclusive \omterm{def:b2:rust-ownership:borrow}{borrows} in use at once: refused
(the program is \code{two_mut.rs}).}\label{out:b2:rust-ownership:two-mut}

\begin{definition}[Borrow checker]\label{def:b2:rust-ownership:checker}
The \emph{borrow checker}\index{borrow checker} is the part of the Rust compiler that checks, for every function,
that each value is used only while owned and not moved, and that the \omterm{def:b2:rust-ownership:borrow}{borrowing rule} holds at every point, using the
types of the function's parameters and results and not the bodies of the functions it calls.
\end{definition}

\begin{omfigure}
\begin{tikzpicture}[font=\scriptsize, v/.style={draw, rounded corners=3pt, fill=omDef!12, minimum width=18mm, minimum height=7mm},
    r/.style={draw, rounded corners=2pt, fill=black!4, minimum width=12mm, minimum height=5mm}]
\node[font=\footnotesize\bfseries] at (1.3,1.9) {allowed};
\node[v] (s) at (0,0) {\texttt{street}};
\node[r] (r1) at (2.6,0.9) {\texttt{\&street}};
\node[r] (r2) at (2.6,0) {\texttt{\&street}};
\node[r] (r3) at (2.6,-0.9) {\texttt{\&street}};
\foreach \r in {r1,r2,r3} \draw[ptr] (\r.west) -- (s.east);
\node[align=center, text=black!60] at (1.3,-1.7) {any number of readers};
\begin{scope}[xshift=4.9cm]
\node[font=\footnotesize\bfseries] at (1.3,1.9) {allowed};
\node[v] (s) at (0,0) {\texttt{street}};
\node[r, fill=omThm!12] (m) at (2.6,0) {\texttt{\&mut street}};
\draw[ptr, thick] (m.west) -- (s.east);
\node[align=center, text=black!60] at (1.3,-1.7) {one writer, alone};
\end{scope}
\begin{scope}[xshift=9.8cm]
\node[font=\footnotesize\bfseries] at (1.3,1.9) {refused};
\node[v] (s) at (0,0) {\texttt{street}};
\node[r, fill=omThm!12] (m) at (2.6,0.5) {\texttt{\&mut street}};
\node[r] (r) at (2.6,-0.5) {\texttt{\&street}};
\draw[ptr, thick] (m.west) -- (s.east);
\draw[ptr] (r.west) -- (s.east);
\draw[red!70!black, very thick] (0.9,-0.95) -- (3.6,0.95) (0.9,0.95) -- (3.6,-0.95);
\node[align=center, text=black!60] at (1.3,-1.7) {a writer and anyone else};
\end{scope}
\end{tikzpicture}
\par\smallskip{\small The \omterm{def:b2:rust-ownership:borrow}{borrowing rule}, for the \omterm{def:b2:rust-ownership:borrow}{borrows} in use at one point of the program. Two writers, like a
writer and a reader, are refused: \cref{out:b2:rust-ownership:two-mut}.}
\end{omfigure}

The rule is what makes the use after free of the opening impossible: dropping \code{name} needs to \omterm{def:b2:rust-ownership:move}{move} it, a \omterm{def:b2:rust-ownership:move}{move}
is a use incompatible with any \omterm{def:b2:rust-ownership:borrow}{borrow} in use, and \code{shown} was in use until the \code{println!}. It also excludes
a whole family of errors that are not about freeing.

% ledger: b2:rustbelt-2018
\begin{proposition}[Aliasing excludes mutation]\label{prop:b2:rust-ownership:no-race}
In a program accepted by the \omterm{def:b2:rust-ownership:checker}{borrow checker}, with no \code{unsafe} code, if at some point a value is changed through
a path (its \omterm{def:b2:rust-ownership:ownership}{owner}, or an exclusive \omterm{def:b2:rust-ownership:borrow}{borrow}), then no other path to that value is in use at that point. In particular,
outside the types built for shared mutation (\code{Cell}, the atomics, \code{Mutex}, all built on \code{UnsafeCell};
\cref{ch:b2:threads-and-synchronisation}), two accesses to one value, one of which is a write, can never be in use at
the same time: in one thread, a value does not change behind a reader; across threads, there is no data race. Those
types change a value through shared references, and still exclude data races.
\end{proposition}

\begin{proof}
A write goes through the \omterm{def:b2:rust-ownership:ownership}{owner}, which requires that no \omterm{def:b2:rust-ownership:borrow}{borrow} be in use, or through an exclusive \omterm{def:b2:rust-ownership:borrow}{borrow}, which by the
\omterm{def:b2:rust-ownership:borrow}{borrowing rule} is the only \omterm{def:b2:rust-ownership:borrow}{borrow} in use; and while an exclusive \omterm{def:b2:rust-ownership:borrow}{borrow} is in use, the \omterm{def:b2:rust-ownership:ownership}{owner} itself cannot be used,
since using it would end or contradict the \omterm{def:b2:rust-ownership:borrow}{borrow}. A second access in use at the same point would be a second
\omterm{def:b2:rust-ownership:borrow}{borrow}, or a use of the \omterm{def:b2:rust-ownership:ownership}{owner}, against the rule. Threads receive values only by \omterm{def:b2:rust-ownership:move}{move} or by \omterm{def:b2:rust-ownership:borrow}{borrow}, so the same
counting applies across them. This argument assumes that the compiler applies the rules correctly to the whole
language and that the library code written with \code{unsafe} respects them. That the rules indeed make the language
safe was proved formally for a realistic subset of Rust in 2018, in a machine-checked proof that also covers
several library types; we admit it.
\end{proof}

\omcode{code/bachelor-2/06-rust-ownership/listings/race.rs}{6}{11}{Two threads incrementing one
counter}\label{lst:b2:rust-ownership:race}

\omoutput[2-16]{code/bachelor-2/06-rust-ownership/outputs/race.txt}{The data race is refused like any second
exclusive \omterm{def:b2:rust-ownership:borrow}{borrow}.}\label{out:b2:rust-ownership:race}

Each closure, \code{|| count += 1}, changes \code{count}, so it holds an exclusive \omterm{def:b2:rust-ownership:borrow}{borrow} of it; the two threads may
run at the same time, so both \omterm{def:b2:rust-ownership:borrow}{borrows} are in use at once. \cref{ch:b2:threads-and-synchronisation} shows the types
that make such a counter possible: an atomic integer, a mutex.

\begin{method}[Fixing a borrow error]\label{met:b2:rust-ownership:fix-borrow}
When the \omterm{def:b2:rust-ownership:checker}{borrow checker} refuses a program, read the three places the message names: where the \omterm{def:b2:rust-ownership:borrow}{borrow} starts, the
conflicting use, and where the \omterm{def:b2:rust-ownership:borrow}{borrow} is used later. Then, in order of preference:
\begin{enumerate}
\item \emph{shorten} the first \omterm{def:b2:rust-ownership:borrow}{borrow}: finish using it before the conflicting use, for instance by copying out the
small value it was needed for (a length, an index);
\item \emph{split} the \omterm{def:b2:rust-ownership:borrow}{borrow}: \omterm{def:b2:rust-ownership:borrow}{borrow} two disjoint parts of the value instead of the whole, two fields, or two halves
of a \omterm{def:b2:rust-ownership:slice}{slice} with \code{split_at_mut};
\item \emph{clone}, if the value is small or the copy is needed anyway, and count what it costs;
\item \emph{restructure}: change who owns what, for instance return an owned value instead of a reference, or keep
indices instead of references.
\end{enumerate}
\end{method}

\section{Lifetimes}\label{sec:b2:rust-ownership:lifetimes}

\begin{definition}[Dangling reference]\label{def:b2:rust-ownership:dangling}
A \emph{dangling reference}\index{dangling reference} is a reference used after the value it refers to was dropped,
or moved: what C calls a dangling pointer, and what the use after free of the opening read.
\end{definition}

\begin{definition}[Lifetime of a reference]\label{def:b2:rust-ownership:lifetime}
The \emph{lifetime of a reference}\index{lifetime (Rust)}, written \code{'a}, is the region of the program in which
the reference may be used. It must lie within the lifetime of the object referred to, the span from its creation
to its \omterm{def:b2:rust-ownership:drop}{drop} or \omterm{def:b2:rust-ownership:move}{move} (the lifetime of an object in Book 3's sense). A function whose result is a reference names, by
lifetime parameters in its signature, the arguments the result may borrow from.
\end{definition}

A function that returns a reference must \omterm{def:b2:rust-ownership:borrow}{borrow} it from one of its arguments: a reference to one of its own locals
would dangle as soon as it returns.

\omoutput[2-8,13-16]{code/bachelor-2/06-rust-ownership/outputs/dangling.txt}{A returned reference with nothing to
\omterm{def:b2:rust-ownership:borrow}{borrow} from (the program is \code{dangling.rs}).}\label{out:b2:rust-ownership:dangling}

When a function takes two references and returns one of them, the compiler cannot guess which; the signature says
it.

\omcode{code/bachelor-2/06-rust-ownership/listings/longest.rs}{2}{29}{A lifetime parameter, and an elided
lifetime}\label{lst:b2:rust-ownership:longest}

\omoutput{code/bachelor-2/06-rust-ownership/outputs/longest.txt}{Both uses stay within the lifetimes of their
arguments.}\label{out:b2:rust-ownership:longest}

\code{fn longest<'a>(x: &'a str, y: &'a str) -> &'a str} reads: for some region \code{'a} in which both arguments
are valid, the result is valid in \code{'a}. A caller may therefore use the result only while both arguments live.
\code{first_letter} needs no annotation: by the \emph{elision rules}, a function with exactly one reference parameter
returns a reference borrowed from it, and a method with \code{&self} returns one borrowed from \code{self}. The caller
checks only the signature, never the body: \code{longest_scope.rs} keeps the result of \code{longest} after \code{b}
is dropped, and is refused even though, with these strings, \code{longest} returns \code{a}.

\omoutput[2-12]{code/bachelor-2/06-rust-ownership/outputs/longest_scope.txt}{The signature promises at most the
shorter lifetime.}\label{out:b2:rust-ownership:longest-scope}

\begin{omfigure}
\begin{tikzpicture}[font=\scriptsize\ttfamily, x=1cm, y=0.5cm]
\foreach \l/\t in {1/let a = String::from(\textquotedbl Orchard Way\textquotedbl);, 2/let result;, 3/\{, 4/{\ \ \ \ let b = String::from(\textquotedbl Mill Row\textquotedbl);},
  5/{\ \ \ \ result = longest(\&a, \&b);}, 6/\}, 7/println!(\textquotedbl\{result\}\textquotedbl);}
  \node[anchor=west] at (0,-\l) {\t};
\draw[very thick, omDef] (6.4,-1) -- (6.4,-7.3);
\node[above, font=\scriptsize\rmfamily, text=omDef] at (6.4,-0.9) {a};
\draw[very thick, omMeth] (6.9,-4) -- (6.9,-6);
\node[above, font=\scriptsize\rmfamily, text=omMeth] at (6.9,-3.9) {b};
\draw[very thick, omThm] (7.4,-5) -- (7.4,-7.3);
\node[above, font=\scriptsize\rmfamily, text=omThm] at (7.6,-4.9) {result};
\draw[red!70!black, thick, dashed] (6.1,-6.3) -- (7.8,-6.3);
\node[font=\scriptsize\rmfamily, text=red!70!black, anchor=west, align=left] at (8.0,-6.3) {b dropped: the result, which may\\ borrow from b, is still to be used};
\end{tikzpicture}
\par\smallskip{\small The lifetimes of \code{longest_scope.rs}, as vertical bars over its source lines. The result
must be valid at the \code{println!}, line 7, but it may borrow from \code{b}, which is dropped at the end of
line 6: no region \code{'a} lies within both, and the program is refused.}
\end{omfigure}

\section{Slices and strings}\label{sec:b2:rust-ownership:slices}

\begin{definition}[Slice, string slice]\label{def:b2:rust-ownership:slice}
A \emph{slice}\index{slice (Rust)} \code{&[T]} is a borrowed view of a contiguous run of values of type \code{T},
inside an array, a vector or another slice: a pointer to its first element and a number of elements. A \emph{string
slice} \code{&str} is a borrowed view of a run of bytes that is valid UTF-8; a \code{String} owns such bytes and can
grow, a \code{&str} only \omterm{def:b2:rust-ownership:borrow}{borrows} them. \code{&v[i..j]} is the slice of the elements from \code{i} up to, but not
including, \code{j}.
\end{definition}

\omcode{code/bachelor-2/06-rust-ownership/listings/slices.rs}{2}{20}{\omterm{def:b2:rust-ownership:slice}{Slices} of an array, a vector and a
string}\label{lst:b2:rust-ownership:slices}

\omoutput{code/bachelor-2/06-rust-ownership/outputs/slices.txt}{One function for every contiguous run of
\code{u32}.}\label{out:b2:rust-ownership:slices}

\begin{omfigure}
\begin{tikzpicture}[font=\scriptsize]
\foreach \v [count=\i from 0] in {65,133,412,98,251} {
  \node[cell, minimum width=11mm, fill=omDef!10] (p\i) at (1.1*\i,0) {\v};
  \node[below=0.5mm of p\i, text=black!60] {\i};
}
\node[left=2mm of p0, font=\scriptsize\ttfamily] {pages};
\node[cell, minimum width=12mm] (sp) at (1.1,2.0) {ptr};
\node[cell, minimum width=12mm, right=0mm of sp] (sl) {len 3};
\node[left=1mm of sp, font=\scriptsize\ttfamily] {middle};
\draw[ptr] (sp.south) -- (p1.north);
\draw[decorate, decoration={brace, mirror, amplitude=4pt}] ($(p1.south west)+(0,-0.45)$) -- ($(p3.south east)+(0,-0.45)$)
  node[midway, below=5pt, font=\scriptsize] {the elements of \texttt{\&pages[1..4]}};
\end{tikzpicture}
\par\smallskip{\small The \omterm{def:b2:rust-ownership:slice}{slice} \code{middle} of \cref{lst:b2:rust-ownership:slices}: two words, a pointer to
element 1 and the length 3, and no copy of the elements. While it is in use, \code{pages} is borrowed: it cannot be
changed or dropped.}
\end{omfigure}

A \code{&str} is a \omterm{def:b2:rust-ownership:slice}{slice} of bytes, but its bounds are byte indices, and UTF-8 encodes é in two bytes: in
the string Café Lane, the bytes 3 and 4 are one character. A \omterm{def:b2:rust-ownership:slice}{slice} that would cut a character in two is not a
string, and the index operation refuses it, at run time when the index is only known then.

\omcode{code/bachelor-2/06-rust-ownership/listings/char_boundary.rs}{8}{10}{A byte index computed at run
time}\label{lst:b2:rust-ownership:char-boundary}

\omoutput{code/bachelor-2/06-rust-ownership/outputs/char_boundary.txt}{A panic, not a broken string: the program
stops (its first lines print a panic without the thread's number, which varies from run to
run).}\label{out:b2:rust-ownership:char-boundary}

A panic, defined in \cref{ch:b2:rust-enums-and-errors}, stops the program in an orderly way: it is a bug reported,
not memory corrupted. The safe ways to cut a string, \code{s.get(0..n)}, which returns \code{None}, or
\code{char_indices}, are the subject of \cref{ch:b2:rust-enums-and-errors}.

\begin{inpractice}[Owned or borrowed?]
Take parameters as \code{&str} and \code{&[T]} rather than \code{&String} and \code{&Vec<T>}: every string and every
vector can be borrowed as a \omterm{def:b2:rust-ownership:slice}{slice}, and the function serves them all. Return owned values (\code{String},
\code{Vec<T>}) when the function creates them, and references only when the result is a part of an argument. Store
owned values in structures, unless the structure is clearly a short-lived view of other data.
\end{inpractice}

\section{Exercises}\label{sec:b2:rust-ownership:exercises}

\begin{exercise}[$\star$]\label{exo:b2:rust-ownership:1}
For each of the five snippets a to e, say whether it compiles, and if not which rule it breaks. Then check with the
session in the solution, which builds each one in turn.
\omcode{code/bachelor-2/06-rust-ownership/listings/exo_compile.rs}{4}{43}{Five snippets}
\end{exercise}

\begin{exercise}[$\star$]\label{exo:b2:rust-ownership:2}
Explain the error with \cref{def:b2:rust-ownership:borrow}, then fix the program in two ways from
\cref{met:b2:rust-ownership:fix-borrow}.
\omcode{code/bachelor-2/06-rust-ownership/listings/exo_e0502.rs}{3}{8}{The shortest name, then one more}
\omoutput[2-10]{code/bachelor-2/06-rust-ownership/outputs/exo_e0502.txt}{The error.}
\end{exercise}

\begin{exercise}[$\star$]\label{exo:b2:rust-ownership:3}
Write \code{fn first_word(s: &str) -> &str}, which returns the part of \code{s} before its first whitespace
character, or all of \code{s}. Why does its signature need no lifetime annotation? Why can the \omterm{def:b2:rust-ownership:slice}{slice} it returns never
cut a character in two? Test it on a word containing é.
\end{exercise}

\begin{exercise}[$\star$]\label{exo:b2:rust-ownership:4}
How many heap allocations does \code{main} make below, and how many times is a buffer enlarged? Count line by line
before reading the output in the solution: an empty \code{Vec} allocates nothing, and the first \code{push} allocates
room for four \code{String}s.
\omcode{code/bachelor-2/06-rust-ownership/listings/exo_clones.rs}{19}{32}{Clones and allocations}
\end{exercise}

\begin{exercise}[$\star\star$]\label{exo:b2:rust-ownership:5}
\code{Noisy} prints its name when it is dropped. Predict the ten lines the program prints, with
\cref{def:b2:rust-ownership:drop}, then check them in the solution. Explain the lines \code{drop c} and
\code{drop a2}.
\omcode{code/bachelor-2/06-rust-ownership/listings/exo_drop.rs}{2}{27}{Who is dropped when?}
\end{exercise}

\begin{exercise}[$\star\star$]\label{exo:b2:rust-ownership:6}
This function returns a reference to its own local. Explain the error, and give the fix the compiler suggests in
\cref{out:b2:rust-ownership:dangling}. What would the equivalent C function return, and when would the bug show?
\omcode{code/bachelor-2/06-rust-ownership/listings/exo_return_local.rs}{3}{6}{Joined names}
\omoutput[2-6]{code/bachelor-2/06-rust-ownership/outputs/exo_return_local.txt}{The error.}
\end{exercise}

\begin{exercise}[$\star\star$]\label{exo:b2:rust-ownership:7}
\code{let (l, r) = (&mut v[..m], &mut v[m..]);} is refused: two exclusive \omterm{def:b2:rust-ownership:borrow}{borrows} of \code{v}. Write \code{add_halves},
which adds each element of the first half of a mutable \omterm{def:b2:rust-ownership:slice}{slice} to the matching element of the second half, with
\code{split_at_mut}. Why may the library give two exclusive \omterm{def:b2:rust-ownership:borrow}{borrows} here, when the compiler would not? Then write
\code{largest_first}, which swaps the largest element to the front, with \code{split_first_mut}.
\end{exercise}

\begin{exercise}[$\star\star$]\label{exo:b2:rust-ownership:8}
Explain why \code{fn f<'a>(x: &'a str, y: &str) -> &'a str} may return \code{x} but not \code{y}. What does the
signature promise the caller, and which call would become unsafe if \code{y} could be returned? Rewrite the
signature so that either may be returned, and say what the caller loses.
\end{exercise}

\begin{exercise}[$\star\star$]\label{exo:b2:rust-ownership:9}
Write \code{fn reverse_words(s: String) -> String}, which reverses the order of the words (separated by single
spaces) in the string's own buffer, with no new allocation: reverse all the bytes, then each word. Prove that the
result is valid UTF-8 even when the words contain multi-byte characters, and test that the returned buffer is the
original one.
\end{exercise}

\begin{exercise}[$\star\star\star$]\label{exo:b2:rust-ownership:10}
Write a generic stack \code{Stack<T>} as a chain of nodes, each a \code{Box<Node<T>>} owning the next node, with
\code{push}, \code{pop}, \code{peek} and \code{peek_mut}. Why does \code{push} need \code{Option::take}? Show that
the default \omterm{def:b2:rust-ownership:drop}{drop} of a chain of a million nodes is a recursion a million calls deep, and write a \code{Drop} that
frees it with a loop. Test the stack against a \code{Vec}.
\end{exercise}

\begin{exercise}[$\star\star\star$]\label{exo:b2:rust-ownership:11}
The loop of \cref{lst:b2:cpp-standard-library:invalidation} read freed memory after a \code{push_back}. (a) Prove,
with \cref{def:b2:rust-ownership:borrow}, that no Rust program without \code{unsafe} can push onto a vector while a
loop iterates over it. (b) Check your argument against the compiler's message for the same loop in Rust, in the
solution. (c) Write the loop the way Rust accepts, and compare it with the fixes of the C++ chapter.
\end{exercise}

\begin{exercise}[$\star\star\star$]\label{exo:b2:rust-ownership:12}
To satisfy the \omterm{def:b2:rust-ownership:checker}{borrow checker}, a port of a C program wrote \code{distinct} as below. The output counts the heap
allocations for $n$ names, each appearing twice. Find the hidden cost, count the allocations exactly as a function of
$n$, and fix the function with no clone of \code{seen}; then give a version in $O(n)$ expected time.
\omcode{code/bachelor-2/06-rust-ownership/listings/exo_quadratic.rs}{13}{23}{Distinct names}
\omoutput{code/bachelor-2/06-rust-ownership/outputs/exo_quadratic.txt}{Allocations for 100, 200 and 400 names.}
\end{exercise}

\section{Problem: Porting the dictionary that crashed}\label{sec:b2:rust-ownership:problem}

\begin{problem}[{Weekend problem --- Porting the dictionary that crashed}]\label{pb:b2:rust-ownership:1}
A word counter in C, a hash table with open addressing in \code{dict.h}, worked for years, then crashed three times,
in three versions, each time after an innocent-looking change. This problem finds the three errors, ports the
counter to Rust, and measures what the port costs.

\needspace{4\baselineskip}\medskip\noindent\textbf{Part I --- Three memory errors in C.}\par\nopagebreak
The function \code{dict_count(d, word)} returns a pointer to the counter of \code{word}, creating it at 0; an
insertion may grow the table, which moves every slot to a new block and frees the old one. Each version below is
built with \code{-DBUG=1}, 2 or 3, and run on the titles of the town's books.
\omcode{code/bachelor-2/06-rust-ownership/listings/dict.c}{9}{36}{\code{dict.c}: the main program and its three
versions}
\omoutput[2-4,6-8,14-18,61-64,66-67,73-75,86-86,88-89,93-93,99-101]{code/bachelor-2/06-rust-ownership/outputs/run_dict_c.txt}{The
key lines of the four runs.}
\begin{enumerate}[beginpenalty=10000]
\item Without bugs, the program prints 34 words; the previous chapter counted 33 in the same titles. Why?
\item Version 1: which pointer is read after its block was freed, and by which call was the block freed?
\item Why does version 1 not crash on every input? Give an input on which it would read nothing freed.
\item Version 2: who owns the string \code{best} points to? Explain the double free.
\item Version 3: \code{dict_count} keeps the caller's pointer instead of a copy. Into which buffer does it point, and
why are some counts already wrong before the crash?
\item Without the sanitizer, what might each version print? Which of the three is the most dangerous?
\end{enumerate}

\needspace{4\baselineskip}\medskip\noindent\textbf{Part II --- An ownership design in Rust.}\par\nopagebreak
\omcode{code/bachelor-2/06-rust-ownership/listings/dict_table.rs}{40}{52}{\code{Dict::count}, from
\code{dict_table.rs}}
\begin{enumerate}[resume, beginpenalty=10000]
\item The slots are \code{Vec<Option<(String, u32)>>}. Who owns each word, and who frees it, and when?
\item \code{grow} replaces the vector with \code{std::mem::replace}, then \omterm{def:b2:rust-ownership:move}{moves} each pair into the new one. Why is no
word copied, and why could \code{grow} not simply read \code{self.slots} while filling a new vector stored in
\code{self.slots}?
\item \code{count} returns \code{&mut u32}. What does its elided lifetime say?
\item Why is there no \code{dict_free} in Rust, and what plays its role?
\item \code{dict.rs} reads each line into a fresh \code{String}. Why can the words of the line be passed to
\code{count} as \code{&str} although the line is dropped at the end of each iteration?
\item Write the order in which the memory of a \code{Dict} holding 34 words is freed at the end of \code{main}.
\end{enumerate}

\needspace{4\baselineskip}\medskip\noindent\textbf{Part III --- The three errors, refused.}\par\nopagebreak
\code{dict_bugs.rs} writes each bug in Rust; the session builds it three times.
\omoutput[2-13,15-19,25-34]{code/bachelor-2/06-rust-ownership/outputs/run_dict_rs.txt}{The program
without bugs, then the three refusals.}
\begin{enumerate}[resume, beginpenalty=10000]
\item Error E0499: relate it to version 1, and to the rule \emph{the pointer is valid until the next insertion}
written in \code{dict.h}'s comment.
\item Error E0507: relate it to version 2. Who may free a \code{String}?
\item Error E0597: relate it to version 3, where a \code{HashMap<&str, u32>} keeps borrowed words.
\item For each of the three, give the fix: the Rust code that does what the programmer meant.
\item Which of the three would also be caught by C++ with \code{std::string} and \code{std::unordered_map}, and which
not?
\item The C comment documented the rule of version 1, and the bug happened anyway. What does the Rust signature of
\code{count} do that the comment did not?
\end{enumerate}

\needspace{4\baselineskip}\medskip\noindent\textbf{Part IV --- The two versions measured.}\par\nopagebreak
The same 4 million words, drawn from 100\,000 by the same generator, are counted by \code{dict.h} and by
\code{dict_table.rs}, three times each, on the author's laptop.
\omoutput{code/bachelor-2/06-rust-ownership/outputs/run_bench.txt}{Three runs of each, on the author's laptop.}
\begin{enumerate}[resume, beginpenalty=10000]
\item Check that the two programs computed the same thing.
\item Compute the mean time of each, and the spread between the runs of one program. Is the difference between the
means larger than the spread?
\item Rust checks the index of every \code{self.slots[i]}. Why does this cost little here?
\item The Rust version makes one heap allocation per new word, the C version one \code{strdup}. Where else could the
two differ in cost?
\item What would a fair comparison need that these runs do not have?
\item State the result of the problem: the three memory errors and what happened to them, and the measured speed of
the two versions.
\end{enumerate}
\end{problem}
